\section{Sophia, the Wisdom of God}

\begin{quotex}
Comte---and for this is due the bulk of his merit and praise---showed more clearly, more decisively, and more fully than all of his predecessors ... showed that the collective whole transcending every single individual by virtue of its internal essence, not merely its external form, and truly fulfilling the individual both ideally and in compete reality, is \emph{humanity} as a living, positive unity that embraces us all. He pointed to the ``Great Being'': \emph{le grand Être}.
\end{quotex}


It is a joy to me when two of my favorite philosophers, \textbf{Vladimir Solovyov} and \textbf{Auguste Comte}, who are as incompatible as unhappy lovers, somehow find a way to make up. Comte, disturbed by the disorder of the French Revolution, sought the antidote. Since Throne and Altar no longer sufficed as the social ordering principle, Comte decided it was up to Science to replace the religion and metaphysic of the Ancien Régime, since Science is inherently conservative. Comte named his system ``Positivism''.

Solovyov had done his master's thesis as a refutation of Positivism! In the meanwhile, Comte realized his system fell short because it did not account for the collective wisdom of mankind. Similarly, Solovyov was developing his own experience of Sophia, or Wisdom. Hence, he returned to the topic of Comte and discovered new riches that were in accord with his own understanding.

\paragraph{Collective Entities}


The existence of collective entities is a metaphysical principle that mostly unknown or certainly poorly understood in the modern West, as \textbf{Rene Guenon} notes in \emph{Man and His Becoming}. As examples, he notes that the six Orthodox Hindu schools darshanas are collective entities. The ``names'' associated with each school therefore are not tied to individuals, but to the collective. He notes:

\begin{quotex}
The formulation of the different darshanas six Hindu schools cannot be related in any way to particular individuals: they are used symbolically to describe what are really intellectual groupings, composed of all those who have devoted themselves to one and the same study over the course of a period the duration of which is no less indeterminable than the date of its beginning.
\end{quotex}


The Ancients used to associate the Lares with geographical locations and the Penates with households. Their worship persisted into the Christian era, but then gave way to a different form. For example, archangels became associated with countries; each village and occupation had its Patron Saint whose days were celebrated with festivities. Generational spirits can affect families. The transition was just nominal. The important point is the knowledge of collective entities, whatever their names, as independent of the individuals of the collectivity.

Solovyov accepts the existence of collectivities and, along with Comte, notes that the individual human being is an abstraction apart from his relationships. Note that Relations are a Category of Being. Solovyov uses geometry to illustrate his meaning: points depend on lines, lines on planes, and planes on bodies. He then states this principle, as applicable to sociology as to geometry:

\begin{quotex}
A whole is primary to its parts and is presupposed by them.
\end{quotex}


He then makes this correspondence:

\begin{enumerate}


\item A sociological point is a single person

\item A line is a family

\item A surface is a nation

\item A three-dimensional figure is a race

\item A completely real body is humanity
\end{enumerate}


In this schema, some points --- or persons --- are more important than others. Just as the centre of a sphere is more significant than other points (because the sphere can be defined by it centre), someone like a Socrates has had worldwide significance, despite his membership in family, nation, etc. A book like \emph{Human Accomplishment} by \textbf{Charles Murray} shows that, in all of history, only 4000 people are of that category. One can quibble with his selection, but no alternative list will be much larger.

Comte and Solovyov are this opposed both to individualism (since the part depends on the whole) and egalitarianism (since not every part is of equal significance).

\paragraph{The Great Being}


The Great Being is therefore humanity in its wholeness. Positive philosophy, starting with sense experience, leads to knowledge of the existence. This Great Being is a female entity, our ``general Mother''. He does not mean a personification, but a person. Solovyov calls her a super-person, since she is not quite a person in the human sense.

While Comte was trying to establish a new religion of humanity, he did not take into account the cult of the Madonna since the Middle Ages. That reached its ``theological apex'' with the decree of the Immaculate Conception. Not only in the West, but a similar feeling was developing in Russia, albeit not yet fully conscious and articulated.


Solovyov references the mysterious icon at the Cathedral of Saint Sophia in Novogord\footnote{\url{https://en.wikipedia.org/wiki/Cathedral\_of\_St.\_Sophia,\_Novgorod}}, whose central figure, the Wisdom of God, seems to be neither the Mother of God nor Christ. Solovyov claims to have unraveled its meaning:

\begin{quotex}
Neither God, nor the Eternal Word of God, nor an angel, nor a holy man, the Great, royal, and feminine Being accepts veneration both from the one who completed the Old Testament and from the foremother of the New Testament. Who can it be other than the truest, purest, and most complete humanity, the highest and all-encompassing form and living soul of nature and the universe, eternally united, and in the process of time uniting with the Divine, and uniting with Him all that is? It is without doubt that this is the full meaning of the Great Being.
\end{quotex}


\paragraph{The Resurrection of the Dead}


The discovery of the Great Being was not Comte's only contribution to Christianity. It is not only the living that comprise Humanity, but also the dead. Solovyov explains the significance:

\begin{quotex}
The dead prevail over the living twofold: as clear role models and as the secret patrons and leaders of the living.
\end{quotex}


Hence, Comte recognized the task of the resurrection of the dead. Solovyov was greatly intrigued by \textbf{Nikolai Fedorov}\footnote{\url{https://www.gornahoor.net/library/cosmism/PhilosophyOfTheCommonTask.pdf}}'s Common Task, which was the beginning of Russian Cosmism. Although Fedorov had in mind the material resurrection of the dead in physical form, Solovyov spiritualized that task. He recognized, in Comte's religion of humanity, the fulfillment of the task.

Solovyov was certainly not a convert to Positivism. Moreover, he rejected Comte's calendar reform whose months were named after prominent figures, with feast days the same. Solovyov obviously preferred the Church calendar with all the saint's days. But this is what we owe to Comte:

\begin{quotex}
He revived in our Christian consciousness ancient and eternal truths, but under new names: the fundamental truth about collective essence of the World Soul---whose simplest name in Christian terms is the Church---and the consummate truth about the life of the dead.
\end{quotex}

Reference:

\textit{The Divine Sophia: The Wisdom Writings of Vladimir Solovyov} by Judith Deutsch Kornblatt


\flrightit{Posted on 2022-06-03 by Cologero}

\begin{center}* * *\end{center}

\begin{footnotesize}
\begin{sffamily}
\texttt{Janet Martha on 2022-06-05 at 15:19 said:}

Thank you for your reflections, Cologero.

I'm reminded of Tomberg's Letter IX on The Hermit. ``The father of empirical science is doubt and its mother is faith... . Newton doubted the traditional theory of gravity, but he believed in the unity of the world and therefore in cosmic analogy. That is why he could arrive at the cosmic law of gravitation in consequence of the fact of an apple falling from a tree. Doubt set his thought in motion; faith rendered it fruitful.''

My imagining of the Great Being changes a lot, and often simply disappears. I found this painting by artist Lizz Daniels, titled ``Sophia Rising'' very evocative. You can see it in the link below.

\url{https://www.academia.edu/49074970/How\_to\_Find\_the\_Virgin\_Sophia\_in\_the\_Foundation\_Stone\_Meditation}


\end{sffamily}
\end{footnotesize}

